\documentclass[11pt,a4paper]{article}
\usepackage[latin1]{inputenc}
\usepackage[T1]{fontenc}
\usepackage{amsmath}
\usepackage{amsfonts}
\usepackage{amssymb}
\usepackage{graphicx}
\usepackage[width=16.00cm, height=24.00cm]{geometry}
\author{Nith Kosal}
\renewcommand{\baselinestretch}{1}
\begin{document}
	\begin{center}
		\large \textbf{Short Biography}
	\end{center}
	
	Nith Kosal is an Young Research Fellow at Future Forum, an independent think tank generating new thinking for a new Cambodia. He holds a double bachelor's degree in Economics and Management from the University of Lumi�re Lyon 2 (France) and the Royal University of Law and Economics (Cambodia) in July 2019, which is the 25th generation. He was a student representative and general coordinator in the French Department of Economics and Management and co-organized the 2019 Eco Study Trip. In his third year, he completed a three-month internship under the Youth Resource Development Program (YRDP) as a program assistant working in support of project planning and coordination of activities and research for good governance in the extractive industries. In 2017, Kosal co-founded the Chea Sim Reak Chey High School Alumni Association and also volunteered in market research with the YRDP and the Koonsoor Kampuchea Training Academy.
	
	In 2018, he was co-author of Reinvigorating Cambodian Agriculture: Transforming from Extensive to Intensive Agriculture, which won the Best Paper Award at the 5th Annual Macroeconomic Conference, organized by the National Bank of Cambodia. He also won first place in the Essay Contest 2019, organized by the Cambodia Development Center on the topic ``How should Cambodia prepare for the Fourth Industrial Revolution?'' Kosal has been a member of the Cambodian Economic Association and has volunteered as a Country Leadership Trainer for Habitat for Humanity, Cambodia since 2018. He has participated in numerous social work projects through various study visits and advocacy activities to promote the protection of natural resources and sustainable development in Cambodia. His research interests include applied economics, macroeconomics, economic development, and international economics.  
	
	\flushright This short biography was last fully updated on \today.
\end{document}